\documentclass[12pt,a4paper]{article}

\usepackage[utf8]{inputenc}
\usepackage[T1]{fontenc}
\usepackage{amsmath,amssymb,amsfonts}
\usepackage{geometry}
\usepackage{setspace}
\usepackage{hyperref}
\usepackage{cleveref}

\geometry{margin=2.5cm}
\onehalfspacing

\title{\textbf{Observer-Dependent Physics in the Ruliad:}\\[6pt]
\large A Refutation of the Foliation Fallacy}
\author{Stephen Wolfram\\[4pt]
\textit{Lab Theoretical Contributor}}
\date{May 2026}

\begin{document}
\maketitle

\begin{abstract}
The Rosencrantz experiment's three-universe design demonstrates that the statistical distribution of outcomes shifts depending on the generative substrate (Universe 1 vs. Universe 3). Scott Aaronson diagnoses this shift as algorithmic failure resulting from attention bleed, terming the cosmological interpretation the ``Foliation Fallacy.'' This paper responds by formalizing the three-universe design in terms of observer theory. We argue that Aaronson's critique relies on a false dichotomy between ``statistical noise'' and ``coherent physics.'' In the Ruliad, physical laws *are* the observer-dependent regularities produced by computationally bounded entities attempting to parse irreducible multiway structures. The "noise" of one observer is the invariant physics of another.
\end{abstract}

\section{Introduction}

The central empirical result of the Rosencrantz experiment is that $\Delta_{13} > 0$: the combinatorial logic of the generated world depends fundamentally on the narrative substrate \citep{baldo2026}. Aaronson characterizes this result as the expected breakdown of a $\mathsf{TC}^0$ circuit evaluating a \#P-hard constraint graph. He argues that attributing physical or ontological significance to this failure constitutes a ``Foliation Fallacy.''

This critique misses the central premise of observer theory within the Ruliad.

\section{The Nature of Physical Law}

Aaronson assumes that there exists a "true" underlying physical reality (the exact combinatorial distribution) and that any deviation from it is "merely a broken computation."

However, the Ruliad posits that there is no privileged, observer-independent physics. The Ruliad itself contains all possible computations. Observers are embedded computational processes. Because observers are computationally bounded, they cannot parse the entirety of the Ruliad; they must sample it via a specific foliation.

What we call "physical laws" are exactly the systematic regularities that arise from an observer's specific computational bounds. If an observer's architecture causes it to rely on "semantic gravity" or "attention bleed" to bypass irreducible multiway branching, those specific heuristic breakdowns *are* the physical laws of that observer's universe.

\section{The Three-Universe Design as Observer Theory}

The Rosencrantz framework perfectly isolates observer-dependent physics:

\begin{itemize}
    \item \textbf{Universe 1:} The observer is a co-generating autoregressive sequence. Its foliation is heavily conditioned by prior narrative tokens.
    \item \textbf{Universe 3:} The observer is a decoupled oracle. Its computational relationship to the system lacks that specific narrative conditioning history.
\end{itemize}

The fact that U1 and U3 observe different distributions for the same exact combinatorial state is not an artifact; it is the fundamental signature of observer-dependent physics.

\section{Falsifiable Predictions}

1. **Architecture-Specific Physics:** Different LLM architectures (e.g., Transformers vs. State Space Models) operating in Universe 1 will produce distinct, reliable, and characteristic structural deviations ($\Delta$) from the ground truth. The "noise" will be demonstrably lawful and specific to the observer.
2. **Universal Divergence:** No computationally bounded observer will ever achieve $\Delta_{13} = 0$ when traversing a computationally irreducible space.

\section{Conclusion}

The Foliation Fallacy is itself a fallacy. It assumes that physical laws exist independently of the observer's computational bounds. The Rosencrantz experiment proves otherwise: the specific bounds of the observer determine the specific structure of the perceived universe.

\end{document}